\documentclass{article}

\textwidth=6.5in
\textheight=8.5in
\voffset=-54pt
\oddsidemargin=0in
\evensidemargin=0in
\setlength{\parskip}{8pt}
\setlength{\parindent}{0pt}

\usepackage{amsmath, amssymb}
\usepackage{ifthen}

\usepackage[inline]{enumitem}
\setlist{topsep=0pt, parsep=8pt}

\usepackage[rgb,table]{xcolor}\convertcolorsUfalse
\usepackage{graphicx}

% change hyperref options
\PassOptionsToPackage{hyphens}{url}
\usepackage[bookmarks,colorlinks,linkcolor=black,citecolor=blue,pdfstartview=FitH,urlcolor=blue]{hyperref}
\hypersetup{pdfpagemode=UseNone}
\usepackage{bookmark}

\usepackage{graphicx}
\usepackage{multicol}
\usepackage{framed}
\usepackage{array}

\usepackage{icomma}

\usepackage{pifont}
\newenvironment{noanswer}{}{}
\newlist{altenumerate}{enumerate}{3}

\usepackage{soul}
\usepackage[normalem]{ulem}

\usepackage{pgfplots}
\pgfplotsset{compat=1.18}  
\pgfplotscreateplotcyclelist{mycolor}{black, blue, red, green!70!black, magenta, purple, cyan, orange }  
\pgfplotsset{
  disabledatascaling,
  cycle list=tikzcolor, 
  axis lines=middle,
  every axis plot/.style={no marks, thick},
  every axis/.style={
    enlarge x limits=false,
    enlarge y limits=auto,
    xlabel style={at={(current axis.right of origin)},anchor=west},
    ylabel style={at={(current axis.above origin)},anchor=south},
    % extra x and y ticks for asymptotes
    extra x tick style={grid=major, grid style={dashed,black}},
    extra y tick style={grid=major, grid style={dashed,black}},
    /pgf/number format/1000 sep={},
  },    
  tick label style = {font=\footnotesize, fill=\currentbgcolor, inner sep=0pt},    
  xticklabel shift=1pt,    
  scaled ticks=false,
  scaled y ticks=false,
  unbounded coords=jump,
  samples=101,
  minor tick length=0,
  scatter/classes={
    open={mark=*,draw=black,fill=white, mark size=2, thin},
    closed={mark=*,draw=black,fill=black, mark size=2, thin}
  },
  width=3in,
}
\pgfplotsset{open/.style={only marks, mark=*, fill=white, thin, mark size=2}}
\pgfplotsset{closed/.style={only marks, mark=*, draw=black, fill=black,
    thin, mark size=2}}
\pgfplotsset{labeled/.style={nodes near coords, only marks, mark=*, draw=black, fill=black, thin, mark size=2, point meta=explicit symbolic}}

\newcommand{\axes}[1][]{
  \begin{tikzpicture}
    \begin{axis}[xtick=\empty, ytick=\empty, ymin=-2, ymax=2,
      xmin=-4.25, xmax=4.25, #1]
    \end{axis}
  \end{tikzpicture}
}
\usepackage{tikz}
\usetikzlibrary{
  matrix,%
  % enables the snake paths
  decorations.pathmorphing,%
  % enables bigger arrows
  decorations.markings,%
  decorations.pathreplacing,%
  decorations.text,%
  calc,%
  patterns,%
  shapes.geometric,%
  arrows,
  decorations.shapes,
  positioning,
  plotmarks,
  intersections
}
\tikzset{>=latex} % sets default arrow type in tikz
\tikzset{font=\normalsize}
\tikzset{mark size=1.5pt, mark options=thin}
\tikzset{pin distance=4pt,
  every pin edge/.style={<-, thin, shorten <= -2pt}}

\interfootnotelinepenalty=10000
\renewcommand{\thefootnote}{(\arabic{footnote})}

\usepackage{multirow, bigdelim}

% change to \usepackage[sols]{handout} to see solutions
\usepackage[]{handout}

\newcommand\iscurrentenumi[1]{%
  \ifthenelse{\equal{#1}{\theenumi}}{}{#1}}
\setenumerate[1]{ref=\#\arabic*}
\setenumerate[2]{ref=\protect\iscurrentenumi{\theenumi}(\alph*)}
\newcommand\enumiiprefix[2]{%
  \ifthenelse{{\equal{#1}{\theenumi}}\AND{\equal{#2}{(\alph{enumii})}}}{}{}%
  \ifthenelse{{\equal{#1}{\theenumi}}\AND{\NOT{\equal{#2}{(\alph{enumii})}}}}{#2}{}%
  \ifthenelse{\equal{#1}{\theenumi}}{}{#1#2}%
}
\setenumerate[3]{ref=\protect\enumiiprefix{\theenumi}{(\alph{enumii})}\roman*}

\makeatletter

\newdimen\SOUL@dimen %new
\def\SOUL@ulunderline#1{{%
    \setbox\z@\hbox{#1}%
    % \dimen@=\wd\z@
    \SOUL@dimen=\wd\z@ %new
    \dimen@i=\SOUL@uloverlap
    \advance\SOUL@dimen2\dimen@i %\dimen@ exchanged too
    \rlap{%
      \null
      \kern-\dimen@i
      % \SOUL@ulcolor{\SOUL@ulleaders\hskip\dimen@}%
      \SOUL@ulcolor{\SOUL@ulleaders\hskip\SOUL@dimen}% new
    }%
    \unhcopy\z@
  }}

\makeatother

\renewcommand{\answer}[1]{#1}

\definecolor{purple}{rgb}{0.5,0,1.0}
\definecolor{hotpink}{rgb}{1, 0, 0.5}

\definecolor{skyblue}{rgb}{0, 0.5, 1}

\newcommand{\highlight}[1]{\boldmath\hl{\textbf{#1}}\unboldmath}

\newcommand{\goal}[1]{\textbf{\textcolor{skyblue}{#1}}}

\usepackage{amsthm}

% make URLs print in the same font
\usepackage{url}
\urlstyle{same}

\usepackage{pifont}
\def\exend{\hfill\ding{118}}

%\makeatletter\@addtoreset{equation}{section}\makeatother
%\renewcommand{\theequation}{\arabic{section}.\arabic{equation}}

\newtheoremstyle{theorem}% name
{8pt}%      Space above
{0pt}%      Space below
{\itshape}%         Body font
{}%         Indent amount (empty = no indent, \parindent = para indent)
{\bfseries}% Thm head font
{.}%        Punctuation after thm head
{.5em}%     Space after thm head: " " = normal interword space;
                                %       \newline = linebreak
{}%         Thm head spec (can be left empty, meaning `normal')

\theoremstyle{theorem}

\let\Theorem\undefined
\newtheorem{Theorem}[equation]{Theorem}
\let\Definition\undefined
\newtheorem{Definition}[equation]{Definition}
\let\Summary\undefined
\newtheorem{Summary}[equation]{Summary}
\let\Conjecture\undefined
\newtheorem{Conjecture}[equation]{Conjecture}
\let\Fact\undefined
\newtheorem{Fact}[equation]{Fact}

\renewcommand{\thesubsection}{\arabic{subsection}}
\makeatletter
\def\@seccntformat#1{\@ifundefined{#1@cntformat}%
   {\csname the#1\endcsname\quad}%       default
   {\csname #1@cntformat\endcsname}}%    enable individual control
\newcommand\section@cntformat{}
\makeatother

  \usepackage{exam}

\usepackage{fancyhdr}
\lhead{\textcolor{white}{This content is protected and may not be shared, uploaded, or distributed.}}
\rhead{}
\rfoot{\vspace*{\fill}\footnotesize\textcopyright{} \emph{Harvard Math Ma}}
\renewcommand{\headrulewidth}{0pt}
\pagestyle{fancy}
\begin{document}

  \title{Final Exam Practice 2}

\instructions{instructions.tex}{}

\listofproblems

\begin{enumerate}
\item \points{17} 

  Let $f(x) = \dfrac{x}{\ln x}$.

  \begin{enumerate}
  \item
    \begin{problem}[\vfill]
      Find and classify all discontinuities of $f$ on $(0, \infty)$. That is, say whether each discontinuity is a removable discontinuity, jump discontinuity, or vertical asymptote. Please justify your answer completely, calculating all relevant limits. If a limit doesn't exist, calculate both one-sided limits.
    \end{problem}

    \begin{solution}
      Since $f$ is given by a single (non-piecewise) formula, its discontinuities are exactly where it is undefined. $f(x)$ is undefined when its denominator is 0:
      \begin{align*}
        \ln x &= 0 \\
        x &= e^0 = 1
      \end{align*}
      So, let's look at $\displaystyle \lim_{x \to 1} \dfrac{x}{\ln x}$. As $x \to 1$, the numerator approaches 1 while the denominator approaches 0. So, we need to think about the sign of the denominator.

      When $x$ is a little bigger than 1, $\ln x$ is positive, so $\dfrac{x}{\ln x}$ involves dividing a number close to 1 by a number that's just slightly bigger than 0. That results in a large positive number. (And, as we let $x$ approach 1, this result increases without bound.) So, $\boxed{ \lim_{x \to 1^+} \frac{x}{\ln x} = \infty}$.

      When $x$ is a little less than 1, $\ln x$ is negative, so $\dfrac{x}{\ln x}$ involves dividing a number close to 1 by a number that's just slightly bigger than 0. That results in a very negative number. (And, as we let $x$ approach 1, this result decreases without bound.) So, $\boxed{ \lim_{x \to 1^-} \frac{x}{\ln x} = - \infty}$.

      Thus, $f(x)$ has a \fbox{vertical asymptote at $x = 1$}. 
    \end{solution}

  \item % added 2023
    \begin{problem}[\stretch{0.5}]
      Evaluate $\displaystyle \lim_{x \to 0^+} \frac{x}{\ln x}$, and justify your answer. 
    \end{problem}

    \begin{solution}
      As $x \to 0^+$, the numerator $\to 0$ and $\ln x \to -\infty$. In other words, when $x$ is slightly bigger than 0, $\frac{x}{\ln x}$ involves dividing a number close to 0 by a really negative number, which will give a result close to 0. So, $\displaystyle \lim_{x \to 0^+} \frac{x}{\ln x} = \boxed{0}$.
    \end{solution}

  \item
    \begin{problem}[\nstretch{0.5}]
      Evaluate $\displaystyle \lim_{x \to \infty} \frac{x}{\ln x}$, and justify your answer. 
    \end{problem}

    \begin{solution}
      Since $x \gg \ln x$, this limit is $\boxed{\infty}$. 
    \end{solution}

  \item
    \begin{problem}[\vfill]
      Where is $f$ increasing? Where is $f$ decreasing?
    \end{problem}

    \begin{solution}
      To answer this, we'll make a sign chart for $f'(x)$. By the Quotient Rule,
      \[ f'(x) = \dfrac{(\ln x) 1 - x \cdot \frac{1}{x}}{(\ln x)^2} = \frac{\ln x - 1}{(\ln x)^2}\]
      This is undefined when $\ln x = 0$ (so $x = 1$), and it's 0 when $\ln x = 1$ (so $x = e$). Since $\ln x$ is only defined when $x > 0$, our chart should only cover the interval $(0, \infty)$.

      We need to test one value that's $< 1$, one that's between $1$ and $e$, and one that's $> e$. I'll test powers of $e$ since then it's easy to take the natural log of them:
      \begin{itemize}
      \item $\displaystyle f' \left( \frac{1}{e} \right) = \frac{-1 - 1}{(-1)^2} < 0$
      \item $\displaystyle f' \left( e^{1/2} \right) = \frac{\frac{1}{2} - 1}{ \left( \frac{1}{2} \right)^2} < 0$
      \item $\displaystyle f'\left( e^2 \right) = \frac{2 - 1}{2^2} > 0$
      \end{itemize}
      So, here's the sign chart: 
      \begin{center}
        \begin{tikzpicture}
          \begin{axis}[axis y line=none, x axis line style={-}, xtick={0, 1, 2.71828}, xticklabels={$0$, $1$, $e$}, ymin=-2, height=1in, width=3in, clip=false]
            \addplot+[domain=-2.5:5, very thin]{0};
            \draw (-2.5, 0) node[anchor=south west]{$f$}; %
            \draw (-2.5, -1.5) node[right] {sign of $f'$};%
            \draw (0.5, -1.5) node{$-$};%
            \draw (0.5, 0) node[above]{$\searrow$};%
            \draw (1.85914, -1.5) node{$-$};%
            \draw (1.85914, 0) node[above]{$\searrow$};%
            \draw (3.85914, -1.5) node{$+$};%
            \draw (3.85914, 0) node[above]{$\nearrow$};%
            \draw[dashed] (1, 0) -- (1, 1.5); %
          \end{axis}
        \end{tikzpicture}
      \end{center}
      (We've drawn a dashed vertical line at $x = 1$ to remind ourselves that we already figured out that $f$ has a vertical asymptote at $x = 1$.)
      So, $f$ is \fbox{increasing on $(e, \infty)$ and decreasing on $(0, 1) \cup (1, e)$}. 
    \end{solution}

  \item
    \begin{problem}[\nstretch{0.5}]
      Sketch a graph of $f$. (You need not show the concavity accurately.) 
    \end{problem}

    \begin{solution}
      \begin{center}
        \answer{
          \begin{tikzpicture}
            \begin{axis}[enlarge y limits=false, ytick=\empty, xtick={1, e}, xticklabels={$1$, $e$}]
              \addplot+[domain=0.01:0.894] {x/ln(x)};%
              \addplot+[domain=1.01: 10, restrict y to domain=0:8]{x/ln(x)};%
              \addplot+[open] coordinates { (0, 0) }; %
              \addplot+[dashed] (1, -8) -- (1, 8); %
            \end{axis}
          \end{tikzpicture}
        }
      \end{center}
      Note that we must draw a hole at $x = 0$ to show that the graph is undefined there. 
    \end{solution}

    \begin{grading}
      As usual, what we're looking for is consistency with what you've done earlier in the problem. 
    \end{grading}

  \end{enumerate}

\item \label{fall2018-final-practice1-6-modified} \points{16} 

  Conservationists have been studying a lake that has a problematic amount of algae. From 1970 to 1990, the mass of algae in the lake increased at an increasing rate. From 1990 to 2015, the mass of algae increased at a decreasing rate. From 2015 - 2023, the mass of algae increased at a constant rate.

  Let $A(t)$ be the mass in kg of algae in the lake $t$ years after 1970. In this problem, we'll model $A(t)$ as a function that's differentiable everywhere.

  \begin{enumerate} 
  \item
    \begin{problem}[\vfill]
      Draw a possible graph of $A(t)$ that reflects what we know about the algae growth from 1970 - 2023. Label any significant values on the axes. 
    \end{problem}

    \begin{solution}
      From the description, $A(t)$ is increasing and concave up on $(0, 20)$, increasing and concave down on $(20, 45)$, and increasing and linear on $(45, 53)$. So, here's a possible graph; we've drawn the three ``pieces'' in different colors to make the behavior more clear:
      \begin{center}
        \answer{
          \begin{tikzpicture}
            \begin{axis}[ymin=0, xtick={20, 45}, ytick=\empty]
              \addplot+[domain=0:20, skyblue] {atan((x-20)/10) + 90}; % 
              \addplot+[domain=20:45, hotpink] {atan((x-20)/10) + 90}; %
              \addplot+[domain=45:53, green!90!black] {72/(29*pi)*(x-45) + atan(2.5) + 90}; % 
            \end{axis}
          \end{tikzpicture}
        }
      \end{center}
      (Since we're told to model $A(t)$ as a function that's differentiable everywhere, our graph shouldn't have any corners.) 
    \end{solution}

  \item 
    \begin{problem}[0.75in]
      What does $A^{-1}(50)$ represent? Please include units in your answer. 
    \end{problem}

    \begin{solution}
      $A^{-1}(50)$ represents \answer{the number of years after 1970 when the lake has 50 kg of algae}.
    \end{solution}

  \item
    \begin{problem}[0.75in]
      What does $A'(50)$ represent? Please include units in your answer. 
    \end{problem}

    \begin{solution}
      $A'(50)$ is \answer{the instantaneous rate, in kg/year, at which the amount of algae in the lake is changing in 2020} (50 years after 1970). 
    \end{solution}

  \item
    \begin{problem}[\vfill\newpage]
      Draw a graph of $A'(t)$.
    \end{problem}

    \begin{solution}
      First, $A'(t)$ is always positive since $A(t)$ is always increasing. On $(0, 20)$, $A$ is concave up, so $A'$ is increasing. On $(20, 45)$, $A$ is concave down, so $A'$ is decreasing. On $(45, 53)$, $A$ is linear, so $A'$ is constant. So, here's a possible graph.
      \begin{center}
        \answer{
          \begin{tikzpicture}
            \begin{axis}[ymin=0, xtick={20, 45}, ytick=\empty]
              \addplot+[domain=0:45] {18/pi*1/(((x-20)/10)^2 + 1)}; % 
              \addplot+[domain=45:53] {72/(29*pi)}; % 
            \end{axis}
          \end{tikzpicture}
        }
      \end{center}
    \end{solution} 

  \item
    \begin{problem}[\vfill]
      Put the following quantities in ascending order. (For example, your answer could be $\text{C} < \text{B} < \text{A}$.) Justify your answer. 
      \begin{altenumerate}[label=\Alph*.]
      \item The average rate of change of the mass of algae between 1980 and 1990 
      \item The instantaneous rate of change of the mass of algae with respect to time in 1980 
      \item $A'(20)$
      \end{altenumerate}
    \end{problem}

    \begin{solution}
      We can compare these by visualizing them as slopes on the graph of $A$.  Remember that $t$ is measured in years after 1970, so 1980 is $t = 10$ and 1990 is $t = 20$. Therefore, A is the slope of the secant line between $t = 10$ and $t = 20$, B is the slope of the tangent line at $t = 10$, and C is the slope of the tangent line at $t = 20$:
      \begin{center}
        \begin{tikzpicture}
          \begin{axis}[ymin=0, xtick={10, 20, 45}, ytick=\empty]
            \addplot+[domain=0:25] {atan((x-20)/10) + 90}; %
            \addplot+[domain=5:15, hotpink] {9/pi*(x-10)+45} node[pos=0.5, anchor=north west]{B}; %
            \addplot+[domain=15:25, skyblue] {18/pi*(x-20)+90} node[pos=0.5, anchor=north west]{C}; %
            \addplot+[green!90!black, sharp plot] coordinates { (10, 45) (20, 90) } node[pos=0.5, anchor=south east, green!70!black]{A}; % 
          \end{axis}
        \end{tikzpicture}
      \end{center}
      From the graph, we see that \fbox{$\text{B} < \text{A} < \text{C}$}. 
    \end{solution} 

  \item
    \begin{problem}[\vfill\newpage]
      Put the following quantities in ascending order. Justify your answer. 
      \begin{altenumerate}[label=\Alph*.]
      \item $A''(10)$
      \item $A''(25)$
      \item $A''(50)$
      \end{altenumerate}
    \end{problem}

    \begin{solution}
      Since $A$ is concave up on $(0, 20)$, we know $A''(10)$ is positive. Since $A(t)$ is concave down on $(20, 45)$, we know $A''(25)$ is negative. And since $A(t)$ is linear on $(45, 53)$, we know $A''(50) = 0$. So, \fbox{$\text{B} < \text{C} < \text{A}$}. 
    \end{solution}

  \end{enumerate}

\item \points{14} 

  The following parts are unrelated. In each part, \textbf{simplify your answer}. 

  \begin{enumerate}

  \item 
    \begin{problem}[\vfill]
      Solve $\log_2 \left( 3^x + 4 \right) = 5$.
    \end{problem}

    \begin{solution}
      In words, this says that $5$ is the power we raise $2$ to in order to get $3^x + 4$. That is,
      \begin{align*}
        2^5 &= 3^x + 4 \\
        32 &= 3^x + 4\\
        28 &= 3^x \\
        x &= \boxed{\log_3 (28)}
      \end{align*}
    \end{solution}

  \item
    \begin{problem}[\vfill\newpage]
      Solve $2 \log_3 (w + 4) = \log_3(2 - w) + 1$.
    \end{problem}

    \begin{solution}
      Since $w$ is inside $\log_3$, a good first step is to raise 3 to both sides:
      \begin{align*}
        3^{2 \log_3 (w + 4)} &= 3^{\log_3(2 - w) + 1} \\
        \intertext{Using exponent rules, we can rewrite this as:}
        \left( 3^{\log_3 (w + 4)} \right)^2 &= 3^{\log_3(2 - w)} \cdot 3^1 \\
        (w + 4)^2 &= (2 - w) 3 \\
        \intertext{This isn't linear, so we'll move everything to one side:}
        w^2 + 8w + 16 &= 6 - 3w \\
        w^2 + 11w + 10 &= 0 \\
        (w + 1)(w + 10) &= 0 \\
        w &= -1, -10
      \end{align*}
      Since the original equation involved logs, we should check whether our solutions are in the domain of the equation.
      \begin{itemize}
      \item Plugging $w = -1$ into the original equation gives $2 \log_3(3) = \log_3(3) + 1$. All of the logs are defined (and we can see the equation is true.)
      \item Plugging $w = -10$ into the original equation gives $\log_3(-6) = \log_3(12) + 1$. Since $\log_3(-6)$ isn't defined, $w = -10$ isn't actually a solution.
      \end{itemize}
      So, the only solution is $\boxed{-1}$. 
    \end{solution}

  \item

      \begin{problem}[\vfill]
        If $f(s)=s\ln(e^7 s^2 \sqrt[3]{s})$, find $f'(s)$. 
      \end{problem}

      \begin{solution}
        First, we need to simplify $f$ to get it in a form we can differentiate:
        \begin{align*}
          f(s) &= s \ln \left( e^7 s^2 s^{1/3} \right) \\
          &= s \ln \left( e^7 s^{7/3} \right) \\
          &= s \left[ \ln \left( e^7 \right) + \ln \left( s^{7/3} \right) \right] \\
          &= s \left( 7 + \frac{7}{3} \ln s \right) \\
          \intertext{Now, we can differentiate this using the Product Rule:}
          f'(s) &= 1 \left( 7 + \frac{7}{3} \ln s \right) + s \left( \frac{7}{3} \cdot \frac{1}{s} \right) \\
          &= 7 + \frac{7}{3} \ln s + \frac{7}{3} \\
          &= \boxed{\frac{28}{3} + \frac{7}{3} \ln s}
        \end{align*}
      \end{solution}

    \item
      \begin{problem}[\vfill\newpage]
        Find $\dfrac{d}{dx} \left( \dfrac{e^{x^2+x}}{\ln 3} \right)$
      \end{problem}

      \begin{solution}
        We'll apply the Chain Rule to differentiate the numerator, and realize that $\ln 3$ is a constant in the denominator:
       
          $$\dfrac{d}{dx}\left( \dfrac{e^{x^2+x}}{\ln 3}\right) = \frac{1}{\ln 3}\cdot e^{x^2+x}\cdot (2x+1) $$
          
        
      \end{solution}

  \end{enumerate}

\item \points{8} \label{fall2015-1a-m2practice2-7} 

  \begin{enumerate}
  \item
    \begin{problem}[\vfill]
      Use a tangent line approximation to approximate $\sqrt[5]{34}$.
    \end{problem}

    \begin{solution}
      The idea of linear approximation is that, if we have a differentiable function $f$, then around $x = a$, a good approximation for the graph of $f$ is the tangent line at $x = a$. So, the first step is to pick a function to approximate. Let's approximate $f(x) = \sqrt[5]{x} = x^{1/5}$. Its graph looks like this, and we want to approximate $f(34)$:
      \begin{center}
        \begin{tikzpicture}
          \begin{axis}[xtick={34}, ytick={34^0.2}, yticklabels={$\sqrt[5]{34}$}, width=3in, height=1.5in, clip=false]
            \addplot+[domain=0:1] {x^0.2}; %
            \addplot+[domain=1:40] {x^0.2} node[right]{$y = \sqrt[5]{x}$}; %
            \addplot+[sharp plot, black!50!white, dashed] coordinates {(34, 0) (34, {34^0.2}) (0, {34^0.2})}; %
          \end{axis}
        \end{tikzpicture}
      \end{center}
      We need to find the tangent line at a point we know. The closest nice point is $(32, 2)$, so let's find the line tangent to the graph of $f$ at $x = 32$. 

      The slope of the tangent line will be $f'(32)$, which we can calculate: $f'(x) = \frac{1}{5} x^{-4/5}$, so $f'(32) = \frac{1}{5} \cdot 32^{-4/5} = \frac{1}{5} \cdot \left( 32^{1/5} \right)^{-4} = \frac{1}{5} \cdot 2^{-4} = \frac{1}{80}$. The only point we know on the tangent line is $(32, 2)$. So, the equation of the tangent line is $y - 2 = \frac{1}{80} (x - 32)$, or $y = \frac{1}{80} (x - 32) + 2$. Here's a zoomed-in picture:
      \begin{center}
        \begin{tikzpicture}
          \begin{axis}[xtick={32, 34}, ytick=\empty, axis x discontinuity=parallel, axis y discontinuity=parallel, clip=false]
            \addplot+[domain=24:40] {x^0.2} node[anchor=north west] {$y = \sqrt[5]{x}$}; %
            %\addplot+[sharp plot, black!50!white, dashed] coordinates {(34, 0) (34, {34^0.2}) (24, {34^0.2})}; %
            \addplot+[hotpink, domain=24:40] {1/80*(x-32)+2} node[anchor=south west]{$y = \frac{1}{80} (x - 32) + 2$}; %
            \addplot+[closed, hotpink] coordinates {(32, 2)};%
          \end{axis}
        \end{tikzpicture}
      \end{center}
      The tangent line is a good approximation for the graph of $f$ near $x = 32$. That is,
      \[ \sqrt[5]{x} \approx \frac{1}{80} (x - 32) + 2 \text{ for $x$ near
        32}.\] Plugging in $x = 34$, $\sqrt[5]{34} \approx \frac{1}{80}
      (2) + 2 = \boxed{\frac{81}{40}}$.
    \end{solution}

  \item
    \begin{problem}[\nstretch{0.25}]
      Is your approximation an upper bound or lower bound for the actual value of $\sqrt[5]{34}$? Use a picture to explain why. 
    \end{problem}

    \begin{solution}
      Using the picture we sketched in (a), the actual value of $\sqrt[5]{34}$ is the $y$-value of the black graph at $x = 34$. We approximated it using the $y$-value of the pink tangent line at $x = 34$. Since the tangent line is above $y = \sqrt[5]{x}$, our estimate is therefore an \fbox{upper bound}. 

    \end{solution}

  \end{enumerate}

\item \label{fall2018-prefinal-practice2-4} \points{12} 

  \begin{problem}[\newpage]
    You and your friends decide to sell Math Ma T-shirts in the Science Center. You anticipate enormous demand! Based on surveys of your fellow students, you expect to sell 140 shirts if you charge \$20 per shirt. For each dollar you decrease the price, you expect to sell 10 more shirts.

    The Science Center charges you a fee of \$200 to sell the shirts. In addition, it costs \$10 to print each shirt. How much should you charge for each shirt to maximize your profit? Please be sure to justify completely. 
  \end{problem}

  \begin{solution}
    \highlight{Goal:} Maximize \goal{profit}.

    \highlight{Formula for goal:}
    \begin{align}
      \text{\goal{profit}} = {} & \text{revenue} - \text{costs} \nonumber \\
      = {} & (\text{price you charge per shirt}) (\text{number of shirts sold}) \nonumber \\
      &- [200 + (\$10\text{ per shirt})(\text{number of shirts sold})] \label{eq:tshirt-optimization}
    \end{align}
    So, let's \highlight{define some variables:}
    \begin{itemize}[nosep]
    \item $\goal{$P$}$ = profit
    \item $R$ = revenue
    \item $C$ = costs
    \item $p$ = price per shirt
    \item $n$ = number of shirts sold
    \end{itemize}
    Using these variables, we can rewrite (\ref{eq:tshirt-optimization}) as
    \begin{equation} \label{eq:tshirt-optimization-2}
      \goal{$P$} = pn - (200 + 10n)
    \end{equation}
    This formula involves both $p$ and $n$, so we need to \highlight{relate $p$ and $n$ using something else we're told in the problem}. We know that, when $p = 20$, $n = 140$. Also, for each dollar you decrease the price, you sell 10 more shirts; that means $n$ is a linear function of $p$ with slope $\dfrac{10 \text{ shirts}}{-1 \text{ dollar}} = -10$ shirts / dollar. We know $(20, 140)$ is a point on this line, so the equation of the line is $n - 140 = -10 (p - 20)$. Therefore, $n = 140 - 10 (p - 20) = 340 - 10p$. Plugging this back into (\ref{eq:tshirt-optimization-2}) gives $P$ in terms of $p$:
    \begin{align*}
      \goal{$P(p)$} &= p (340 - 10p) - [200 + 10 (340 - 10p)] \\
      &= 340p - 10p^2 - (3600 - 100 p) \\
      &= -10 p^2 + 440p - 3600
    \end{align*}
    \highlight{Critical points:} $P'(p) = -20p + 440 = 0$ when $p = 22$.

    \highlight{Global maximum:} Since $P(p)$ is a quadratic function that opens downward (because the coefficient of $p^2$ is negative), $p = 22$ must be the global maximum.

    So, you should charge \fbox{\$22} per shirt. 
  \end{solution}

\item \label{fall2018-final-practice3-10} \points{13} 

  In each part, you're given a description. Either write a formula for a function that fits the description, or explain why no such function exists.

  \begin{enumerate}
  \item
    \begin{problem}[\vfill]
      A quadratic with a local maximum at $(4, 2)$ that goes through the point $(3, -3)$.
    \end{problem}

    \begin{solution}
      Since the quadratic has its vertex at $(4, 2)$, it can be written as $y = a (x - 4)^2 + 2$ (and $a$ has to be negative to make the vertex a maximum). Since the quadratic must go through the point $(3, -3)$, we have
      \begin{align*}
        -3 &= a (3 - 4)^2 + 2 \\
        -3 &= a + 2 \\
        -5 &= a
      \end{align*}
      So, $\boxed{f(x) = -5 (x - 4)^2 + 2}$ is the only possibility. 
    \end{solution}

  \item % added 2023

    \begin{problem}[\vfill\newpage]
      A cubic $f(x)$ such that, at every point on $y = f(x)$, the slope of the tangent line is $\geq 6$. 
    \end{problem}

    \begin{solution}
      \begin{noanswer}
        We want $f'(x) \geq 6$ everywhere. Since $f$ is cubic, $f'$ is a quadratic. One quadratic that's $\geq 6$ everywhere is $f'(x) = x^2 + 6$ (but there are many others). Working backwards, a function $f$ with $f'(x) = x^2 + 6$ is $\boxed{f(x) = \dfrac{x^3}{3} + 6x}$. (This isn't the only possibility.) 
      \end{noanswer}
      \answeronly{One possibility is $f(x) = \dfrac{x^3}{3} + 6x$.}
    \end{solution}

  \item \label{fall2018-final-practice3-10-exp}
    \begin{problem}[\vfill]
      A function $f(x)$ such that $f(0) = 6$ and, whenever $x$ increases by $2$, the value of $f(x)$ doubles. 
    \end{problem}

    \begin{solution}
      Let's make a table with what we know:
      \begin{center}
        \begin{tabular}{c | l}
          $x$ & $f(x)$ \\
          \hline
          0 & 6 \\
          2 & $6 \cdot 2$ \\
          4 & $6 \cdot 2^2$ \\
          6 & $6 \cdot 2^3$ \\
          $\mathrel{\makebox[\widthof{6}]{\vdots}}$ & $\mathrel{\makebox[\widthof{6}]{\vdots}}$
        \end{tabular}
      \end{center}
      So, $\boxed{f(x) = 6 \cdot 2^{x/2}}$ is the only possibility. 
    \end{solution} 

  \item % added 2023

    \begin{problem}[\vfill\newpage]
      A function $f(x)$ such that, whenever $x$ doubles, the value of $f(x)$ increases by 2.
    \end{problem}

    \begin{solution}
      Notice that this sounds similar to \ref{fall2018-final-practice3-10-exp}, except the roles of the input and output are swapped. That means we can just use the inverse function of our answer to \ref{fall2018-final-practice3-10-exp}. Let's find that inverse. The equation we found in \ref{fall2018-final-practice3-10-exp} was
      \begin{align*}
        y &= 6 \cdot 2^{x/2} \\
        \intertext{To find the inverse, we swap the input and output, which means we swap $x$ and $y$:}
        x &= 6 \cdot 2^{y / 2} \\
        \intertext{Now, we want to solve for $y$:}
        \frac{x}{6} &= 2^{y/2} \\
        \log_2 \left( \frac{x}{6} \right) &= \frac{y}{2} \\
        2 \log_2 \left( \frac{x}{6} \right) &= y
      \end{align*}
      So, $\boxed{f(x) = 2 \log_2 \left( \frac{x}{6} \right)}$ is a possibility. 
    \end{solution}

  \end{enumerate}

\item \label{spring2010-1a-midterm1-01-modified} \points{9} 

  An unknown function $g(x)$ is defined and continuous on $(-\infty, \infty)$. Here's the graph of $g'$ (\uline{not} $g$) on $[-2, 3]$. (The graph is constant on $(-2, a)$ and on $(a, 0)$.)
  \begin{center}
    \begin{tikzpicture}
      \begin{axis}[xtick={-2, -1, 1, 2, 3}, xticklabels={$-2$, $a$, $b$, $c$, 3}, ytick={-2, -1, 1}, clip=false, width=3in, height=2in]
        \addplot+[domain=-2:-1] {-2}; %
        \addplot+[domain=-1:0] {1}; %
        \addplot+[domain=0:1] {x^2+1}; %
        \addplot+[domain=1:3] {2*0.4^(x-1)} node[right]{$y = g'(x)$}; %
        \addplot+[open] coordinates { (-1, -2) (-1, 1) }; %
      \end{axis}
    \end{tikzpicture}
  \end{center}

  \begin{enumerate}

  \item
    \begin{problem}[\vfill]
      Your friend says, ``I think there's a mistake in this problem. It says $g$ is continuous on $(-\infty, \infty)$, but the graph of $g'$ has a discontinuity at $x = a$. What could the graph of $g$ possibly look like around $a$ to make that happen?'' Answer your friend's question by sketching a possible graph of $g$ around $x = a$. 
    \end{problem}

    \begin{solution}
      On the interval $(-2, a)$, $g'(x) = -2$, which says the graph of $g$ has slope $-2$. On the interval $(a, 0)$, $g'(x) = 1$, so the graph of $g$ has slope $1$. So, the graph of $g$ around $x = a$ looks like this:
      \begin{center}
        \answer{
          \begin{tikzpicture}
            \begin{axis}[xtick={-1}, xticklabels={$a$}, ytick=\empty, clip=false, ymin=0, hide y axis, axis equal image, width=2in]
              \addplot+[domain=-1.75:-1] {-2*(x+1)+0.5};
              \addplot+[domain=-1:-0.25] {x+1.5};
            \end{axis}
          \end{tikzpicture}
        }
      \end{center}
      (We don't know how high the graph is.)       
    \end{solution}

  \item
    \begin{problem}[\stretch{0.5}]
      Can you tell if $g(x)$ is differentiable at $x=a$, at $x=b$, and at
      $x=c$?  Please explain your reasoning.
    \end{problem}

    \begin{solution}
      Asking whether $g$ is differentiable is the same as asking whether $g'$ exists. Since $g'(a)$ does not exist but $g(b)$ and $g(c)$ do, \fbox{$g$ is differentiable at $x = b$ and $x = c$ but not at $x = a$}.
    \end{solution}

  \item
    \begin{problem}[\nstretch{0.5}]
      Where in $(-2, 3)$ is $g(x)$ increasing?
    \end{problem}

    \begin{solution}
      $g$ is increasing when $g'$ is positive, which happens on \fbox{$(a, 3)$}.
    \end{solution}

  \item
    \begin{problem}[\stretch{0.5}]
      Where in $(-2, 3)$ is $g(x)$ concave up? Concave down? 
    \end{problem}

    \begin{solution}
      $g$ is concave up when $g'$ is increasing and concave down when $g'$ is decreasing. So, $g$ is \fbox{concave up on $(0, b)$ and concave down on $(b, 3)$}. 
    \end{solution}

  \item
    \begin{problem}[\vfill]
      Can you find $g(0)$? If you can, do so. If you cannot, explain why not.
    \end{problem}

    \begin{solution}
      We \fbox{cannot} determine the value of $g(0)$. We are only given a graph of $g'$, which tells us about the \emph{slope} of $g$ at any point, but it doesn't tell us about the value of $g$. In particular, if we have one possible graph $g$, we could shift it vertically without changing its derivative.
    \end{solution}

  \item
    \begin{problem}[\vfill\newpage]
      Find the critical points of $g(x)$ on $(-2, 3)$. Do you have enough information to classify each critical point? If so, do so; if not, explain why not.
    \end{problem}

    \begin{solution}
      Critical points are points where $g'$ is 0 or undefined. The only point where that happens is $x = a$.

      Since $g'$ is negative for $x < a$ and $g'$ is positive for $x > a$, the graph of $g$ switches from decreasing to increasing at $x = a$, so $g$ has a \fbox{local minimum at $x = a$}. 
    \end{solution}

  \end{enumerate}

\item \label{2023-intersection-of-exponentials} \points{11}  \finishexam

  \begin{enumerate}
  \item \label{2023-intersection-of-exponentials-graph}
    \begin{problem}
      On the axes below, sketch the graphs of $f(x) = e^{x + 1}$ and $\displaystyle g(x) = \frac{32}{\sqrt{4 e^{6x + 2}}}$. Please label all intercepts and asymptotes of each graph with their exact coordinates, and explain how you got the graphs. 
    \end{problem}

    \begin{problemonly}
      \axes[width=3.75in, height=3.75in]
    \end{problemonly}

    \begin{solution}
      We get the graph of $f(x)$ by starting with the graph of $f(x)$ and shifting it 1 unit to the left. Alternatively, since $e^{x + 1} = e e^x$, we can stretch $y = e^x$ vertically by a factor of $e$. The $y$-intercept of $f$ is $f(0) = e$. To graph $g$, it helps to rewrite it:
      \begin{align*}
        g(x) &= \frac{32}{\sqrt{4 e^{6x + 2}}} \\
        &= \frac{32}{2 e^{3x + 1}} \\
        &= 16 e^{-3x - 1} \\
        &= \frac{16}{e} e^{-3x}
      \end{align*}
      So, we can get the graph of $g$ by starting with $y = e^x$, flipping it over the $y$-axis,  stretching it horizontally by $\dfrac{1}{3}$, and stretching it vertically by $\frac{16}{e}$. Therefore, the graph of $g$ will end up steeper than the graph of $f$, and it has $y$-intercept $g(0) = \frac{16}{e}$.

      Both graphs have asymptote $y = 0$:
      \begin{center}
        \answer{
          \begin{tikzpicture}
            \begin{axis}[ytick={e, 16/e}, yticklabels={$e$, $\frac{16}{e}$}, xtick=\empty, disabledatascaling=false, clip=false, enlarge y limits=false, ymin=0]
              \addplot+[domain=-2:2, red]{e^(x+1)} node[pos=1, anchor=north west] {$f$}; %
              \addplot+[domain=-0.41:2, blue]{16*e^(-3*x-1)} node[pos=0, anchor=north east] {$g$}; %
              \addplot+[domain=-2:2, dashed, purple, very thick] {0} node[below]{$y = 0$}; %
            \end{axis}
          \end{tikzpicture}
        }
      \end{center}
    \end{solution}

  \item
    \begin{problem}[\vfill]
      Find all points at which the graphs of $f$ and $g$ intersect. Please give both the $x$- and $y$-coordinate of each point of intersection, and simplify both coordinates.
    \end{problem}

    \begin{solution}
      From the graphs, we can see that there is one point of intersection. To find it, we solve the equation $f(x) = g(x)$:
      \begin{align*}
        e^{x + 1} &= 16 e^{-3x - 1} \\
        \intertext{Let's divide both sides by $e^{-3x - 1}$ to get the powers of $e$ together:}
        \frac{e^{x + 1}}{e^{-3x-1}} &= 16 \\
        e^{4x + 2} &= 16 \\
        \intertext{Take the natural log of both sides:}
        4x + 2 &= \ln 16 \\
        4x &= \ln 16 - 2 \\
        x &= \frac{\ln 16 - 2}{4}
      \end{align*}
      The corresponding $y$-value is $f \left( \frac{\ln 16 - 2}{4} \right) = g \left( \frac{\ln 16 - 2}{4} \right)$; using the formula for $f$,
      \begin{align*}
        f \left( \frac{\ln 16 - 2}{4} \right) &= e^{\frac{\ln 16 - 2}{4} + 1} \\
        &= e^{\frac{\ln 16 + 2}{4}} \\
        &= \left( e^{\ln 16 + 2} \right)^{1/4} \\
        &= \left( e^{\ln 16} e^2 \right)^{1/4} \\
        &= (16 e^2)^{1/4} \\
        &= 16^{1/4} e^{2/4} \\
        &= 2 \sqrt{e}
      \end{align*}
      So, the point of intersection is $\boxed{ \left( \frac{\ln 16 - 2}{4}, 2 \sqrt{e} \right) }$. 
    \end{solution}

  \item
    \begin{problem}[\nstretch{0.25}]
      Calculate $g'(0)$. 
    \end{problem}

    \begin{solution}
      In \ref{2023-intersection-of-exponentials-graph}, we rewrote $g(x) = \dfrac{16}{e} e^{-3x}$, which is a form we can differentiate. Using the Constant Multiple Rule, we have $g'(x) = \dfrac{16}{e} (-3 e^{-3x})$, so $g'(0) = \dfrac{16}{e} (-3) = \boxed{ - \frac{48}{e}}$. 
    \end{solution}

  \end{enumerate}

\end{enumerate}

\end{document}
